\section{Introduction}\label{sec:intro}
\addcontentsline{toc}{part}{\protect\tocpart{}{}{Introduction}}

Let $X$ be a smooth complex projective variety. The Hodge conjecture asserts
that the cycle class map
\[
  \cl^{p}_{X}\colon \CH^{p}(X)_{\QQ}\longrightarrow \Hdg^{p}(X)
  \coloneqq H^{2p}(X,\QQ)\cap H^{p,p}(X)
\]
is surjective for every $p$. For $p=1$ this is the theorem of Lefschetz on
$(1,1)$ classes.%
\footnote{The rational coefficients are essential. The theorem of Lefschetz
holds integrally in degree two, but the integral analogue fails in higher
degree: Atiyah and Hirzebruch found torsion classes of type $(p,p)$ that are
not classes of cycles, and Koll\'ar found non-torsion ones. Every statement in
this paper concerns $\QQ$-coefficients.}
Beyond degree two the best understood Hodge classes are those of abelian
varieties, and the first of them that divisors do not produce are the Weil
classes. Weil exhibited them on abelian varieties with an action of an
imaginary quadratic field whose two eigenspaces on the tangent space have the
same dimension, showed that they are Hodge classes outside the subring
generated by divisor classes, and proposed them as test cases for the Hodge
conjecture \cite{Weil77}. The construction works for every CM field $F$ in
place of the quadratic one. On a general abelian variety of Weil type the
Hodge ring is generated by divisor classes together with the Weil classes, a
flat $F$-line of classes of middle degree, so for these varieties the Hodge
conjecture is the algebraicity of the Weil classes.

Markman proved that the Weil classes are algebraic on the abelian sixfolds of
split Weil type for every imaginary quadratic field, and, through a product
argument of Schoen \cite[Proposition~10]{Sch98}, on every abelian fourfold of
Weil type \cite{Mar25a,Mar25b}.%
\footnote{In the normalisation of \Cref{def:disc} the discriminant is the
class of $(-1)^{n}\det H$ in $\QQ^{\times}_{>0}/\Nm_{K/\QQ}(K^{\times})$.
Markman uses the class of $\det H$ itself, in which the same sixfolds have
discriminant $(-1)^{3}=-1$.}
With the description of the Hodge rings of abelian varieties of dimension at
most five by Moonen and Zarhin \cite{MZ99}, this gives the Hodge conjecture
for abelian fourfolds and fivefolds. His sheaves are built from secant sheaves
on an abelian variety of half the dimension, and they extend along the family
because they satisfy an equivariant form of semiregularity. This paper is
about what can be proved for every CM field and every dimension at once.

\subsection{The main theorem}

Fix a CM field $F$ of degree $2m$, an integer $n\ge2$ and a discriminant
$\delta$. The abelian varieties of $(F,n,\delta)$-Weil type are the members of
one family over a connected period domain $\cD_{F}$ of dimension $mn^{2}$, a
product of $m$ copies of the bounded symmetric domain of $\SU(n,n)$, and their
Weil classes form a flat local system of rank $2m$ over it
(\Cref{setup:cmfield}, \Cref{prop:cmweil}). We call the set
\[
  \Sigma_{F}=\{\,s\in\cD_{F}:\text{the Weil classes of }A_{s}\text{ are
  algebraic}\,\}
\]
the \emph{algebraic locus} of the family. A property holds for the
\emph{very general member} if it holds outside a countable union of proper
closed analytic subsets of $\cD_{F}$.%
\footnote{A subset of $\cD_{F}$ is \emph{meagre} if it is contained in a
countable union of closed subsets with empty interior. Since $\cD_{F}$ is a
complex manifold, the Baire category theorem shows that a meagre set has
dense complement, so a statement that holds off a meagre set holds on a dense
set of members; the very general members are those off a particular meagre
set, a countable union of proper analytic subsets.}

\begin{theorem}[Main theorem]\label{thm:main}
For every CM field $F$, every $n\ge2$ and every discriminant $\delta$, the
algebraic locus $\Sigma_{F}$ of the Weil classes has the following
properties.
\begin{enumerate}[label=\textup{(\roman*)},leftmargin=2.6em,itemsep=2pt]
\item It contains an explicit point $s_{0}$ at which $A_{s_{0}}$ is isogenous
  to a product of abelian varieties with complex multiplication by $F$, and
  at which the Weil classes are polynomials with rational coefficients in
  divisor classes.
\item It is dense in $\cD_{F}$.%
\footnote{Density comes from the Hecke orbit of $s_{0}$: the group
$G_{F}(\QQ)$ of rational points of the unitary group of the hermitian form
acts on $\cD_{F}$ through isogenies that respect the $F$-action and the
polarisation up to a scalar, so it preserves the algebraicity of the Weil
classes, and the closure of $G_{F}(\QQ)$ contains the identity component of
$G_{F}(\RR)$, which acts transitively on $\cD_{F}$; the proof uses the
rational isotropic vectors of the hermitian form, which exist by the Hasse
principle.}
\item It is the preimage of a countable union of Zariski closed subsets of the
  Shimura variety of the family, and either $\Sigma_{F}=\cD_{F}$ or
  $\Sigma_{F}$ is meagre.
\item The following are equivalent: $\Sigma_{F}$ is closed; $\Sigma_{F}$ has
  an interior point; the Weil classes are algebraic on every member of the
  family; the Hodge conjecture holds for the very general member of the
  family.
\end{enumerate}
\end{theorem}

For $n=1$ the Weil classes have degree two and are algebraic by the theorem
of Lefschetz, and all four statements hold trivially. The proof of
\Cref{thm:main} is completed in \Cref{sec:cmfields}, after the imaginary
quadratic case has been treated in full in \Cref{part:basepoints}.
\Cref{fig:mainlocus} draws the picture it describes: a connected domain, a
dense set built from one explicit point, and a decomposition of that set into
countably many closed algebraic strata.

The theorem carries the Hodge conjecture for the very general abelian variety
of Weil type, for every CM field, every dimension and every discriminant, to
a single topological statement about one explicit set: the dense set
$\Sigma_{F}$, which is a countable union of closed algebraic subsets of a
connected Shimura variety, has to be closed. The rest of the paper is
organised around that statement. \Cref{part:obstructions} shows why the
classical constructions do not produce Weil classes and fixes the explicit
form of the classes; \Cref{part:basepoints} proves the theorem for imaginary
quadratic fields; \Cref{part:semiregularity} determines what closing
$\Sigma_{F}$ requires of a single object at a single point;
\Cref{part:cmfields} completes the proof for every CM field and computes the
requirement exactly for quartic fields; and \Cref{part:beyond} places the
Weil classes inside the whole conjecture (\Cref{fig:parts}).

The paper also proves the Hodge conjecture itself on a family of varieties
of every dimension. Let $X\subset\PP^{N}$ be a complete intersection of $c$
diagonal hypersurfaces $\sum_{i}a_{ki}x_{i}^{d}=0$ of degree $d$. It is
\emph{of Vandermonde type} if $a_{ki}=\mu_{i}\sum_{l=1}^{c}g_{kl}\lambda_{i}^{l-1}$
for distinct numbers $\lambda_{0},\dots,\lambda_{N}$, nonzero numbers
$\mu_{i}$ and an invertible matrix $(g_{kl})$; for $c\ge2$ this says that the
columns of $(a_{ki})$ are distinct points of a rational normal curve, and
every smooth complete intersection of two diagonal hypersurfaces of the same
degree is of this type. Such an $X$ is isomorphic to a variety
$X_{d,N-c}(\lambda)$ that depends only on $\lambda$ (\Cref{thm:vandermonde}),
and \emph{very general} refers to $\lambda$. For $d=2$ we call $X$ a
\emph{Vandermonde quadric intersection}.

\begin{theorem}[Diagonal complete intersections]\label{thm:introvandermonde}
Every rational Hodge class is algebraic on the very general Vandermonde
quadric intersection, in every dimension and in every projective space. The
same holds for the very general complete intersection of Vandermonde type of
diagonal cubics or quartics of dimension at most seven, or of diagonal
sextics of dimension at most five, and for every smooth complete
intersection of two diagonal cubics, of two diagonal quartics or of two
diagonal sextics in $\PP^{6}$.
\end{theorem}

The first two statements are \Cref{thm:vgvandermonde}(iii) and the third is
\Cref{cor:twodiagonal}(ii). The variety $X$ is a quotient of a power of a
generalised Fermat curve, so its Hodge classes come from abelian subvarieties
of the Jacobian of that curve. For quadrics the monodromy of the curve along
a chain of vanishing cycles is the full symplectic group on every eigenspace,
so every Hodge class comes from powers of polarisations and no theorem on
abelian varieties of higher dimension is needed. For cubics, quartics and
sextics we compute the monodromy by an induction on the number of branch
points (\Cref{prop:cyclicmonodromy}), and the classes it leaves are divisor
classes, Weil classes of abelian fourfolds of Weil type, and Weil classes of
abelian sixfolds of split Weil type, all algebraic by the theorems of
Lefschetz and Markman \cite{Mar25a,Mar25b}. In odd dimension every Hodge
class is a multiple of a power of the hyperplane class. So the bound seven
for cubics and quartics is the point where abelian varieties of dimension
eight appear, and the bound five for sextics comes from abelian sixfolds of
non-split Weil type, which occur in dimension six. Taking the theorems of
the preprints discussed below as hypotheses, these bounds rise to nine and
seven (\Cref{thm:weilsix}), and the Hodge conjecture holds on every power of
the very general member in every degree (\Cref{cor:vgalldegrees}); for one
or two equations that last statement is a theorem of \cite{OAI26a}, proved
there independently.

\begin{figure}[htbp]
\centering
\includegraphics[width=\textwidth]{fig_roadmap}
\caption{The five parts of the paper. Parts I to IV lead to
\Cref{thm:main}; all five lead to the closure theorem of
\Cref{sec:final}, by which the Hodge conjecture is equivalent to the two
statements (F2) and (F3$'$), both of which remain open.}
\label{fig:parts}
\end{figure}

\begin{figure}[htbp]
\centering
\includegraphics[width=\textwidth]{fig_mainlocus}
\caption{The algebraic locus of \Cref{thm:main}, drawn over a real slice of the
period domain $\cD_{F}$. Left: the explicit CM base point $s_{0}$ of part (i);
its Hecke orbit $G_{F}(\QQ)\cdot s_{0}$, which is dense by part (ii); three of
the strata $\Sigma_{\underline{P},\underline{q}}$ of part (iii), each closed
and algebraic, through points of the orbit; the Hodge loci $Z_{t}$, dashed,
one of them through $s_{0}$, off which the Hodge group of the member is all of
$G_{F}$ (\Cref{lem:cmgeneric}); and a very general point $s$. Right: the two
cases of the dichotomy of part (iii). By part (iv) the Hodge conjecture holds
for the very general member exactly in the upper case, where one stratum is
the whole domain.}
\label{fig:mainlocus}
\end{figure}

\subsection{The proof}

The first step is a base point. In \Cref{thm:cmbasepoint} the member $A_{s_{0}}$
is isogenous to $B_{\Phi}^{\,n}\times B_{\bbar\Phi}^{\,n}$, where $B_{\Phi}$ is
an abelian $m$-fold with complex multiplication by $F$ of a CM type determined
by the polarisation, and the Weil classes there are written out as
polynomials with rational coefficients in $F$-balanced divisor classes. For
an imaginary quadratic field $K$ there is a second base point in every family,
with quaternionic multiplication: whenever an algebra $(-d,b)_{\QQ}$ with
$b>0$ acts compatibly with $K$, the Weil line is spanned by the real and
imaginary parts of $\Theta^{n}$ for an explicit $(1,1)$ class $\Theta$, and
the discriminant of the family is the class of $b^{n}$, which can be any
discriminant (\Cref{thm:divcriterion}, \Cref{thm:quatweil},
\Cref{thm:everyfamily}). Products of abelian surfaces of Weil type give a
third, since the Weil class is multiplicative:
$\omega(A_{1})\,\omega(A_{2})=\omega(A_{1}\times A_{2})$ for the
$K$-valued Weil class $\omega$ (\Cref{thm:weilmult}, \Cref{thm:everyn}). None
of these base points uses anything beyond the theorem of Lefschetz on $(1,1)$
classes.

The second step moves the base point. An element of $G_{F}(\QQ)$, where
$G_{F}=\Res_{F_{0}/\QQ}\SU(V,H)$ and $F_{0}$ is the totally real subfield,
induces after clearing denominators an isogeny between two members and fixes
the Weil classes, so $\Sigma_{F}$ is a union of $G_{F}(\QQ)$-orbits
(\Cref{lem:Ginvariance}). The hermitian form is isotropic by the Hasse
principle, the rational points of the unipotent radicals of the parabolic
subgroups are dense in their real points, and these real points generate the
identity component of $G_{F}(\RR)$; so every orbit is dense
(\Cref{prop:densealg}, \Cref{thm:cmpropagation}(ii)).

The third step is properness. The relative Hilbert scheme of the universal
family, for one tuple of Hilbert polynomials, is proper over the Shimura
variety,%
\footnote{By Grothendieck's construction the Hilbert scheme of a projective
morphism, for a fixed Hilbert polynomial, is projective over the base; here
the base is the Shimura variety of the family, after a level structure that
makes the universal family exist, and the morphism is the universal abelian
scheme.}
and the class of a flat family of cycles is locally constant, so the
set of points at which a Weil class is represented by cycles with prescribed
Hilbert data is Zariski closed; $\Sigma_{F}$ is the countable union of these
strata (\Cref{thm:closedlocus}). The Baire category theorem then gives the
dichotomy of part (iii) (\Cref{cor:allornothing}). Part (iv) needs one more
fact, that the very general member has Hodge group $G_{F}$; we prove it
directly, from the simplicity of the real factors $\SU(n,n)$
(\Cref{lem:cmgeneric}, \Cref{cor:verygeneral}).

For an imaginary quadratic field everything can be written in coordinates.
On the standard model of \Cref{sec:explicit}, built from a $K$-vector space of
dimension $2n$ with a hermitian form of signature $(n,n)$, the Weil line has
the integral basis
\[
  \omega_{1}=\sum_{|T|\ \mathrm{even}}(-1)^{|T|/2}d^{\,n-|T|/2}\,w^{(T)},
  \qquad
  \omega_{2}=\sum_{|T|\ \mathrm{odd}}(-1)^{(|T|-1)/2}d^{\,n-(|T|+1)/2}\,w^{(T)},
\]
where $T$ runs over the subsets of $\{1,\dots,2n\}$ and $w^{(T)}$ is the
monomial taking $y_{j}$ for $j\in T$ and $x_{j}$ otherwise
(\Cref{thm:explicitgens}). A member with Hodge group $\SU(V,H)$ has exactly
three Hodge classes in degree $2n$ and one in every other degree
(\Cref{thm:hdgcount}), so the Hodge conjecture for it, in every degree at
once, is the algebraicity of the one class $\omega_{1}$, with $2^{2n-1}$
integer coefficients (\Cref{cor:whatisleft}). By Deligne's theorem
\cite{Del82} that class is absolutely Hodge at every point,%
\footnote{A class is absolutely Hodge if its image under every automorphism
$\sigma$ of $\CC$, through the comparison of Betti cohomology with
algebraic de Rham cohomology of the conjugate variety $X^{\sigma}$, is again
a Hodge class. Algebraic classes have this property, so it is a necessary
condition for algebraicity; Deligne proved it for every Hodge class on an
abelian variety.}
and the Hodge
conjecture for the family is equivalent to its algebraicity at every point of
one connected domain (\Cref{cor:exactform}).

\subsection{Closing the algebraic locus}

\enlargethispage{2pt}By \Cref{thm:main}(iv) the Hodge conjecture for the very general member of a
Weil family is the closedness of $\Sigma_{F}$, and by (iii) it is enough to
show that one stratum has an interior point. Two statements would do this,
and the paper determines what each of them asks.

The first bounds the Hilbert data. If at every point of the orbit of the
base point the Weil class is represented by a cycle whose Hilbert data comes
from one finite set, then finitely many strata cover a dense set and one of
them is everything (\Cref{thm:degreebound}). Read on an arbitrary
representative this bound is equivalent to its conclusion
(\Cref{thm:p1equivalent}); read on the transported cycle, the one
representative that can be named, it fails, because the transport multiplies
the degree of a component by $c^{2n}$ and divides it only by the order of a
stabiliser (\Cref{thm:p1forces}). So the bound restates the problem rather
than reducing it.
The bound is needed only on a small set: by the theorem of Pila, Shankar
and Tsimerman on the Andr\'e--Oort conjecture \cite{PST},%
\footnote{In the form used here: an irreducible closed subvariety of a
Shimura variety of abelian type that contains a Zariski dense set of CM
points is a special subvariety. Hence a sequence of CM points that leaves
every proper special subvariety is Zariski dense, and so is every infinite
subsequence of it.}
it suffices along
one sequence of CM points that leaves every proper special subvariety, and
such a sequence exists inside the orbit (\Cref{thm:escape}). Every
member at which this paper exhibits algebraic Weil classes has endomorphisms
beyond the field, so it lies in the locus of members with smaller Hodge
group. That locus contains the base point, is dense, is stable under
$G_{F}(\QQ)$ and is a countable union of closed algebraic strata, exactly as
$\Sigma_{F}$ is, and it is meagre (\Cref{prop:speciallocus},
\Cref{cor:speciallocus}). The properties listed in \Cref{thm:main} therefore
do not decide the dichotomy by themselves, and a proof that $\Sigma_{F}$ is
everything has to reach a member with full Hodge group.

The second is semiregularity. By the theorem of Buchweitz and Flenner, in the
form of Pridham for perfect complexes \cite{BF03,Pri24},%
\footnote{The semiregularity map of a perfect complex $E$ on $X$ sends
$\Ext^{2}(E,E)$ to $\bigoplus_{q}H^{q+2}(X,\Omega^{q}_{X})$ by the trace of
the product with powers of the Atiyah class of $E$; it carries the
obstruction to deforming $E$ along a first order deformation of $X$ to the
obstruction to keeping $\ch(E)$ of Hodge type. If it is injective, $E$
deforms, formally and then along the whole family, in every direction in
which $\ch(E)$ stays a Hodge class. Bloch introduced it for local complete
intersection subvarieties.}
a perfect complex at
one base point whose Chern character is a nonzero multiple of a Weil class
plus a polynomial in the polarisation, and whose semiregularity map is
injective, extends along the whole family and makes $\Sigma_{F}$ everything
(\Cref{thm:semiregclosure}, \Cref{thm:cmpropagation}(v)). The paper computes
the injectivity condition exactly. The annihilator of the Weil class in
$H^{0,2}$ is carried by the polarisation isomorphically onto the tangent space
of the family (\Cref{lem:annihilator}, \Cref{thm:anntangent}), so no object
whose Chern character lies on the Weil line and on which $H^{2}(\cO_{A})$ acts
faithfully is semiregular (\Cref{thm:rankzero}). When the Chern character is
exactly a Weil class the criterion is never met, for any CM field and any
$n\ge2$: a semiregular object would extend to the very general Weil torus of
the field, a complex torus that is not algebraic, on which every coherent
sheaf has vanishing Chern character strictly between degree zero and the top
degree (\Cref{prop:weiltorussheaves}, \Cref{thm:p2false},
\Cref{thm:cmpurefalse}). With powers of the polarisation added the criterion
becomes one number: for a Chern character $N\omega+\sum_{k}c_{k}\theta^{k}$
and $n\ge3$ every object has $\dim\Ext^{2}(E,E)\ge(4+\rho)n^{2}-2n$, where
$\rho\le3$ is the rank of the Hankel matrix of the numbers $k!\,c_{k}$, and
an object attaining this bound meets the criterion
(\Cref{thm:p2primenumber}).%
\footnote{That is, $\rho$ is the rank of the $(2n-1)\times3$ matrix with
entries $\mu_{i+j}$, $0\le j\le2n-2$, $0\le i\le2$, where $\mu_{k}=k!\,c_{k}$
(\Cref{setup:p2prime}). A vector in its kernel is a linear recurrence of order
at most two among the $\mu_{k}$, and $\rho$ does not change when the character
is multiplied by $e^{t\theta}$.}
A flatness condition on the deformations of the object suffices in place of
injectivity (\Cref{prop:p2flat}). By descent and scalar extension the
statement is needed only for the split families of the CM fields of degree at
least four (\Cref{prop:descent}), and for one family it is equivalent to the
Lefschetz standard conjecture for the total space of the family over a curve
through a base point (\Cref{cor:weilequiv}).

For a quartic CM field and $n=2$ we compute the rank exactly for every
shape. It takes fifteen values; the least, $80$, is taken only by the pure
shape, which the Weil tori exclude, so an object meeting the criterion has
$\dim\Ext^{2}(E,E)\ge88$ and a polynomial part that is exponential or linear
at each real place (\Cref{thm:quarticrank}, \Cref{cor:quarticleast}). The
Hodge locus of every such character is determined: it is the Weil tori of the
field only for the characters of $\theta^{4}$ shape, which are therefore
excluded at rank $88$, and the polarised family itself for every other
character of that rank (\Cref{prop:quarticlocus}, \Cref{cor:quarticlocus}).

\subsection{Secant objects}

The objects that a secant construction supplies are examined in
\Cref{sec:secantobjects}. On the split member $X\times\widehat{X}$ of
\Cref{sec:construction} the secant plane is written down in every dimension,
and a secant sheaf has rank $a$, $c_{1}=bt$, $c_{2}=\tfrac{ad+b^{2}}{2}t^{2}$
and Bogomolov discriminant $\Nm_{K/\QQ}(b+a\sqrt{-d})\,t^{2}$
(\Cref{thm:plane}, \Cref{thm:secantexists}). No object whose Chern character
is an effective combination of exponentials of divisor classes contributes to
the Weil line, by a positivity argument with the norm form
(\Cref{thm:positivity}), and no object with a line bundle summand off the
polarisation line meets the criterion, so no sum of line bundles does in any
dimension (\Cref{thm:linebundle}, \Cref{cor:noline}). Convolutions of line
bundles, their iterated cones, are cut down by the maps between the pieces.
Fewer than $2n$ pieces with nondegenerate differences never meet the
criterion (\Cref{thm:fewpieces}). On $E_{0}^{2n}$, $E_{0}=\CC/\ZZ[i]$, the
signed Chern characters of $2\cdot4^{n-1}$ line bundles add up to a real
Weil class (\Cref{prop:weilpieces}), and with three multiples of the
polarisation no convolution of these pieces meets the criterion at $n=3$
(\Cref{thm:noconvolution}), nor at $n=4$ when the shifts of the pieces lie
in six consecutive values (\Cref{cor:efoursix}). The weakened criterion
of \cite[Question~11.4]{Mar25b} reduces, for the transform of
$\cF\boxtimes\cF$, to the single identity
$\mathrm{At}(\cF)^{2}=-d\Theta^{2}\cdot\id$ for the sheaf $\cF$ on the
fourfold (\Cref{thm:factor}); a local computation at an isolated double point
of the support decides it in the negative for Markman's candidate in
dimension eight (\Cref{thm:candidatefails}), and at a quartic CM field a
curve on which a complex is locally free is tangent to at most one of the two
eigenplanes of the real multiplication, which decides Markman's explicit pair
\cite[Example~11.2.7]{Mar25c} in the same way (\Cref{thm:quarticlocal},
\Cref{cor:markmanquarticfails}).

The supports are constrained as well. A Cohen--Macaulay support carries no
local obstruction, while two smooth branches of codimension two meeting
transversally at a point carry one (\Cref{cor:cmdichotomy}). For a smooth support, a surface $S$ in the
fourfold, Grothendieck--Riemann--Roch fixes every numerical invariant, and the
Bogomolov--Miyaoka--Yau inequality gives $d\le b^{2}$ for the twist
$b\Theta$, which leaves only $d\in\{1,3,5,7\}$ for the twist $3\Theta$ of
Markman's candidate (\Cref{thm:smoothinvariants},
\Cref{cor:fourdiscriminants}). Complete
intersections are excluded for every $b$ and $d$ (\Cref{thm:noci}), two-term
complexes of powers of the polarisation at every rank, twist and
discriminant (\Cref{thm:nosplit}), and a Hilbert--Burch resolution of a
secant support has rank at least two (\Cref{prop:lowrankburch}); the secant
class is an integer combination of classes of line bundles exactly when
$d\equiv15$ or $23\pmod{24}$ (\Cref{thm:lattice}).

\subsection{The classical constructions}

\Cref{part:obstructions} explains why two classical lines of attack do not
produce Weil classes, and the explanations are different in kind. The first
imports an object from a projective space, restricts it to the dual abelian
variety and transports it back by the Fourier--Mukai transform. Its Chern
character stays in the subring generated by the polarisation
(\Cref{thm:divisor}), because the cohomological transform carries powers of
the dual polarisation to rational multiples of powers of the polarisation in
closed form (\Cref{thm:fmtheta}); and a construction equivariant for the
field lands in the isotypic component of the norm character, which the Weil
line misses, since it sits at the two ends of the grading
(\Cref{thm:character}). The second is induction on dimension through
hyperplane sections, and the restriction map kills every primitive class of
middle degree (\Cref{thm:hyperplane}); below the middle degree the return
trip is the Lefschetz standard conjecture (\Cref{thm:returntrip}). The
Kuga--Satake construction, which changes the variety rather than cutting it,
reaches the Weil line only in dimension two (\Cref{prop:normpart}).%
\footnote{The construction attaches to a polarised Hodge structure $T$ of
weight two with $h^{2,0}=1$ an abelian variety $A$ with
$H^{1}(A,\QQ)=C^{+}(T)$, the even Clifford algebra, and an embedding of
$T$ into $\End(H^{1}(A,\QQ))$ by Hodge structures, which is a Hodge class on
the product of a variety carrying $T$ with $A\times A$; the Kuga--Satake
Hodge conjecture asks that this class be algebraic.} A rational
class of the Weil line has a $\Kx$-orbit spanning a plane, so any
construction reaching it must be bilinear in a pair of inputs exchanged by
$\Gal(K/\QQ)$ (\Cref{prop:orbitplane}), and Markman's construction is of that
kind.

\subsection{Weil classes and the Hodge conjecture}

\Cref{part:beyond} places the Weil classes inside the conjecture. Two
statements carry all of it. The first, (F2), asks that every Hodge class on an
abelian variety outside the subring generated by divisor classes and Weil
classes of quotients be algebraic; multiplying by an exceptional class on the
square of a Mumford fourfold moves every Hodge class of every abelian variety
outside that subring, so (F2) is equivalent to the Hodge conjecture for all
abelian varieties (\Cref{prop:f2isab}). The second, (F3$'$), asks that every
Hodge class on every smooth projective variety be algebraic modulo the images
of Hodge classes of abelian varieties under algebraic correspondences
(\Cref{def:f3prime}). The Hodge conjecture is equivalent to (F2) and (F3$'$)
together (\Cref{prop:f3prime}, \Cref{thm:final}).

We prove (F3$'$) for every variety whose cohomology is reached from abelian
varieties by correspondences, a class closed under products, surjective
images, blow-ups, projective bundles and Hilbert schemes of points on
surfaces, which contains every Fermat hypersurface and every K3 surface whose
Kuga--Satake correspondence is algebraic (\Cref{prop:f3prime},
\Cref{prop:aclosure}); for every variety of dimension at most three, and on
a variety of dimension $k$ in the degrees $0$, $2$, $2k-2$ and $2k$; for
every fourfold and fivefold whose zero-cycles are supported in dimension at
most three, among them every uniruled fourfold and every rationally connected
fivefold (\Cref{prop:f3primesmall}, \Cref{prop:f3primefibration}); for every
fourfold and fivefold dominated by a variety of that class
(\Cref{prop:f3primedominant}); for every smooth ample divisor of odd
dimension in a variety of that class, and in even dimension outside the
middle degree and for a very general member (\Cref{prop:f3primeample});
and, by a toric compactification whose boundary is again of the same kind, for
every hypersurface of simplex type, so for every smooth Delsarte hypersurface
of every dimension and every cyclic cover branched along one, whose Hodge
conjecture then follows from that of Fermat varieties of one degree
(\Cref{thm:simplextype}, \Cref{cor:delsarteall}). A complete intersection
of diagonal hypersurfaces whose coefficients are of Vandermonde type, which
holds for every smooth complete intersection of two of them, is a quotient
of a power of a generalised Fermat curve, so it satisfies (F3$'$), and its
Hodge conjecture reduces to abelian varieties attached to the characters of
the curve; for two diagonal hypersurfaces of degree $3$, $4$ or $6$ in
$\PP^{6}$ these have dimension at most five, and the Hodge conjecture holds,
the exceptional classes being Weil classes of abelian fourfolds of Weil type
(\Cref{thm:vandermonde}, \Cref{cor:twodiagonal}). At a very general
configuration the monodromy of the curve leaves only the Hodge classes that
come from polarisations and from Weil classes; we compute that monodromy for
cyclic covers of degree three, four and six by an induction on the number of
branch points (\Cref{prop:cyclicmonodromy}), so the Hodge conjecture holds
for the very general such intersection of quadrics in every dimension, of
cubics and quartics up to dimension seven, and of sextics up to dimension
five, in every projective space (\Cref{thm:vgvandermonde}). On the powers, Hilbert
schemes and moduli spaces of a K3 surface whose transcendental field is a CM
field, and on the fourfolds of K3$^{[2]}$ type with field $\QQ$ or a CM
field, the Hodge conjecture itself holds (\Cref{prop:k3powers},
\Cref{prop:k3ntype}).

The first instance of (F2) is a pair of classes on the square of a Mumford
fourfold, and \Cref{sec:beyond} studies it in detail. The two classes are
the real multiplication by a totally real cubic field on the primitive part of
$H^{2}$, and they are algebraic once the Kuga--Satake class of one K3 surface
of Picard number thirteen is algebraic (\Cref{prop:mumfordrm},
\Cref{thm:mumfordks}); they are algebraic at every CM point
(\Cref{thm:mumfordcm}); they are rigid, their Hodge locus being the compact
Mumford curve (\Cref{thm:mumfordrigid}); and they are represented by
algebraic cycles at every closed point of every reduction of the curve modulo
a prime, so that their algebraicity on every member is a bound on the Hilbert
data of those cycles (\Cref{prop:modp}, \Cref{thm:modpdensity}). For an
abelian scheme over a curve whose invariant classes are algebraic the
Lefschetz standard conjecture holds on the total space
(\Cref{thm:lefschetzfamily}), and over a compact Shimura curve the total space
of a Mumford family satisfies the Hodge conjecture
(\Cref{prop:mumfordfivefold}).

The structure that \Cref{thm:main} finds in one family is present in every
family that carries (F2) or (F3$'$) (\Cref{sec:families}). Every Hodge class
of every abelian variety lies in an abelian scheme with a member of CM type,
and (F2) is equivalent to the statement that in every such scheme the
algebraic locus of a flat Hodge class that is algebraic at a CM member has an
interior point (\Cref{prop:cmfamily}, \Cref{thm:f2propagation}); the base
points of the Weil families are the seeds, through Andr\'e's theorem on the
Hodge classes of abelian varieties of CM type. (F3$'$) is equivalent to the
absence of Hodge classes off the part of the cohomology reached from abelian
varieties by algebraic correspondences (\Cref{prop:abpart}). On the squares of
K3 surfaces with real multiplication, where (F3$'$) is first open, the
algebraic locus of the real multiplication contains every CM point, is dense,
is stable under rational Hodge isometries, and is either the whole family or
meagre (\Cref{thm:rmlocus}). In each case the property still missing is the
interior point (\Cref{thm:twostand}).

Seven preprints released by OpenAI in September and October 2026
\cite{OAI26a,OAI26b,OAI26c,OAI26d,OAI26e,OAI26f,OAI26g} announce the Hodge
conjecture for abelian varieties of CM type, the algebraicity of the Weil
classes on the split abelian eightfolds of imaginary quadratic fields and of
the Kuga--Satake class of every K3 surface, the Hodge conjecture for products
of K3 surfaces, and the Hodge conjecture for the powers of Hodge generic
abelian covers of curves and of diagonal hypersurfaces. They have not been
refereed. \Cref{sec:preprints} takes their theorems as hypotheses and
derives, with complete proofs, what they give here: the Weil classes of every
abelian sixfold of Weil type, by a product argument that lowers the dimension
(\Cref{thm:weilsix}); the Hodge conjecture on every power of the
Kuga--Satake variety of a K3 surface, through a spanning lemma for Clifford
algebras, and so on every power of a Mumford fourfold (\Cref{thm:kspowers},
\Cref{cor:mumfordall}); every K3 surface in the class of varieties reached
from abelian varieties, and the
fourfolds of K3$^{[2]}$ type whose N\'eron--Severi space is isotropic or
represents $-2$, as when the Picard number is at least four
(\Cref{cor:k3all}, \Cref{prop:k3twosmall}); the Hodge conjecture
on every hypersurface of simplex type and on every power of the very general
diagonal complete intersection of Vandermonde type (\Cref{cor:simplexhc},
\Cref{cor:vgalldegrees}); and the CM points inside every algebraic locus
(\Cref{cor:f2cm}). No other section depends on them. Neither (F2) nor
(F3$'$) follows: the Weil classes of eightfolds of nontrivial discriminant
and of quartic CM fields, and (F3$'$) beyond the varieties reached from
abelian varieties, remain open (\Cref{rem:remains}).

Finally, the implications of the paper and of the literature form a finite
rule set, and its consequences can be determined by hand. They give the
minimal sets of statements from which the conjecture follows
(\Cref{thm:closure}, \Cref{rem:f2rule}): (F3$'$) together with the Lefschetz standard
conjecture, with the variational Hodge conjecture for algebraic classes, or
with (F2), besides two sets that are each equivalent to the conjecture.%
\footnote{The Lefschetz standard conjecture asks that the inverse of the
hard Lefschetz isomorphism $L^{\dim X-k}\colon H^{k}(X,\QQ)\to
H^{2\dim X-k}(X,\QQ)$ be induced by an algebraic correspondence. The
variational Hodge conjecture asks that a flat section of Hodge classes over
the base of a smooth projective family that is algebraic at one fibre be
algebraic at every fibre.}
\Cref{sec:final} collects the results in one theorem.

Every proof is written to be checked by hand. Where an argument rests on a
finite calculation, the calculation is given where it is used, as an explicit
identity, a determinant or a case analysis. The archive \cite{BMB26} holds
programs that repeat much of this arithmetic and a Lean~4 certificate of part
of it; no statement of the paper rests on them.
